% summary.tex
% Вставь этот файл вместо текущего topics/summary.tex

\section*{Convexity }
\textbf{Def.} \(f\) convex \(\iff\) epi\((f)\) convex \(\iff f(\theta
x+(1-\theta)y)\le\theta f(x)+(1-\theta)f(y)\) \(\forall\theta\in[0,1]\).

\textbf{Characterizations (use these):}
\begin{itemize}[leftmargin=*,itemsep=-1pt]
  \item Twice diff.: \(\nabla^2 f(x)\succeq0\ \forall x\) (evals for
    \(x^\top Qx\)).
  \item 1st-order: \(f(y)\ge f(x)+\nabla f(x)^\top(y-x)\ \forall x,y\).
  \item Monotone grad: \((\nabla f(x)-\nabla f(y))^\top(x-y)\ge0\).
  \item \(f^{**}=f\) (proper, closed epi).
  \item Restriction to any line is convex.
\end{itemize}

\textbf{Preserving ops:} nonnegative lin.\ comb., pointwise
\textbf{max}, affine comp.\(f(Ax+b)\), intersection of convex sets.

\textbf{Does \emph{NOT} imply convex (traps + counterex.):}
\begin{itemize}[leftmargin=*,itemsep=-1pt]
  \item Unique/global or single local min (\(f(x)=|x|\) on
    \([-1,1]\), \(1\) outside).
  \item \(\nabla^2 f\) bounded below (must be \(\succeq0\)).
  \item \(|\nabla f(x)-\nabla f(y)|\ge\mu|x-y|\) (\(\mu>0\)).
  \item Pointwise \textbf{min} of convex: e.g.\(f_1=(x+1)^2\), \(f_2=(x-1)^2\).
  \item Product \(f_1f_2\): e.g.\(x\cdot(-x)=-x^2\) (concave).
  \item \(\exp(f(x))\) or \(g(x)=1/\ln x\) (\(x>1\)) (Hessian neg.).
  \item \(\sqrt{x}\) on \([0,1]\): no finite Lip.\ const.\ (\(\nabla
    f\to\infty\) at 0).
\end{itemize}

\textbf{Examples:}
\begin{itemize}[leftmargin=*,itemsep=-1pt]
  \item \(\frac12 x^2 + p\ln x\) on \([1,\infty)\): convex \(\iff
    p\le1\) (\(\nabla^2=1-p/x^2\succeq0\)).
  \item \(1/x_i\) convex on \(\{x>0\}\).
  \item Quadratic \(x^\top Qx\) convex \(\iff Q\succeq0\).
\end{itemize}
\textbf{Tactic (MC/TF):} Test Hessian/1st-order or give counterex.
Local min = global \(\nRightarrow\) convex.

\section*{Conjugate \(f^*(y)=\sup_x(y^\top x-f(x))\)}
\(f\) proper convex closed: \\
\( f^{**}=f\). \(y\in\partial
f(x)\iff x\in\partial f^*(y)\).

\textbf{4-Step Recipe:}
\begin{enumerate}[leftmargin=*,itemsep=-1pt]
  \item Set \(\nabla_x(y^\top x - f(x))=0\implies\nabla f(x)=y\).
  \item Solve for stationary \(x_y\).
  \item If \(x_y\in\text{int}(\text{dom }f)\), \(f^*(y)=y^\top x_y - f(x_y)\).
  \item Restricted domain \([a,b]\): evaluate boundaries (\(x_y<a\to
    x=a\), \(x_y>b\to x=b\)).
\end{enumerate}

\textbf{Quiz Examples:}
\begin{itemize}[leftmargin=*,itemsep=-1pt]
  \item \(f(x)=-x^2\), \(x\in[-1,1]\): \(f^*(y)=1+|y|\) (V-shape
    plot, label axes).
  \item \(f(x)=x^2+x\), \(x\in[-1,1]\): stationary \(x=\frac12(y-1)\),
    \[
      f^*(y)=
      \begin{cases}
        -y & y\le-1\\
        \frac14(y-1)^2 & -1\le y\le3\\
        y-2 & y\ge3
      \end{cases}
    \]
  \item \(f(x)=-\ln x\) (\(x>0\)): \(f^*(y)=-1-\ln(-y)\) (\(y<0\)),
    \(\infty\) else.
  \item \(f(x)=e^x\): \(f^*(y)=y\ln y-y\) (\(y>0\)), \(\infty\) else.
  \item \(f(x)=\frac12 x^\top Qx\): \(f^*(y)=\frac12 y^\top Q^{-1}y\).
\end{itemize}
\textbf{Tactic:} Plot if asked (label axes). Stationary often
quadratic; always check boundaries.

\section*{Subdifferential \& Optimality}
\(\partial f(\bar x)=\{v\mid f(x)\ge f(\bar x)+v^\top(x-\bar x)\ \forall x\}\).

\textbf{Optimality:} \(0\in\partial f(x^*)\iff x^*\) global minimizer
(unique if strict convex).

\textbf{Example \(f(x)=|x|\):}
\[
  \partial f(x)=
  \begin{cases}
    \{-1\} & x<0\\
    [-1,1] & x=0\\
    \{1\} & x>0
  \end{cases}
\]
Minimizer at points where subdifferential plot contains \(0\)
(vertical segment over \(y=0\)).

\textbf{Tactic (plot Q):} Find point where subdiff.\ interval/segment
includes \(0\).

\textbf{Projections (PGD):} \(x_{k+1}=\mathrm{prox}_{T,C}(x_k-T\nabla f(x_k))\).
\begin{itemize}[leftmargin=*,itemsep=-1pt]
  \item \(\|x\|_\infty\le1\): clamp each coord to \([-1,1]\).
  \item \(\|x\|_2\le1\): \(z/\max(1,\|z\|_2)\).
  \item \(\ell_1\)-ball: soft-thresholding.
\end{itemize}

\section*{Duality, Lagrangian, KKT}

\begin{sstTitleBox}{
    Lagrange Duality
  }
  \begin{centering}
    \vspace{-1.5mm}
    \begin{equation}\text{Consider }
      f^\star = \inf_{x\in\mathcal{\mathbb{R}}^n}f(x)
      \text{ s.t.}\ g(x)\le0,h(x)=0
      \label{eq:dual}
    \end{equation}

    \begin{sstFullFrame}
      {\color{white}
        \vspace{-2mm}
        \[
          \begin{aligned}
            & \textbf{Lagrangian}
            & \hspace{-4mm}  \mathcal{L}(x,\lambda,\nu)
            & = f(x) + \lambda^T g(x)+\nu^T h(x)
            \\
            \\
            & \textbf{Dual Function}
            & \hspace{-4mm}    d(\lambda,\nu)            & = \inf_{x
            \in \mathcal{\mathbb{R}}^n}\mathcal{L}(x,\lambda,\nu)
        \end{aligned}\]
        \vspace{-2.7mm}
      }
    \end{sstFullFrame}
  \end{centering}
\end{sstTitleBox}

\textbf{Weak duality:} \(d^*\le p^*\) \textbf{always}.

\textbf{Strong duality + KKT} (convex + Slater: \(\exists\hat x\),
\(g_i(\hat x)<0\)):
\begin{align*}
  \nabla_x\mathcal L(x^*,\lambda^*,\nu^*)&=0,\\
  g(x^*)&\le0,\ h(x^*)=0,\ \lambda^*\ge0,\\
  \lambda^{*\top}g(x^*)&=0\ (\text{compl.\ slackness}).
\end{align*}

\begin{sstTitleBox}[BrickRed]{\textbf{\large
      KKT  Conditions
    }
  \normalsize(Karush-Kuhn-Tucker)}
  \begin{sstOnlyFrame}[BrickRed]
    If Slater's condition holds and
    OP is convex
    $\rightarrow$
    $x^\star \in \mathbb{R}^{n}$ is a minimizer of the primal
    and
    $(\lambda^\star\ge0,\nu^\star)\in\mathbb{R}^{n_g}\times\mathbb{R}^{n_h}$
    is a maximizer of the dual $\Leftrightarrow$
    to the following statements:
  \end{sstOnlyFrame}
  \begin{sstFullFrame}[BrickRed]
    \vspace{-3mm}
    \color{white}
    \small
    \[
      \begin{aligned}
        \textbf{KKT-1 } & \text{\footnotesize (Stationary
        Lagrangian)} \hspace{-2mm} &  &
        \nabla_x\mathcal{L}(x^\star,\lambda^\star,\nu^\star)  =  0
        \\
        \textbf{KKT-2 } &
        \text{\footnotesize (primal feasibility)}
        && g(x^\star)\le0, h(x^\star)= 0
        \\
        \textbf{KKT-3 } & \text{\footnotesize (dual feasibility)}
        &  & \lambda^\star  \ge 0,\quad \nu^\star \in \mathbb{R}^{n_h}
        \\
        \textbf{KKT-4 } & \text{\footnotesize (complementary}
          &  & \lambda^{\star T} g(x^\star)
          =                          0
          \\
        & \text{\footnotesize slackness)}
        &  & \nu^{\star T} h(x^\star)                              =  0
    \end{aligned}\]
    \vspace{-4mm}
  \end{sstFullFrame}
  \begin{sstOnlyFrame}[BrickRed]
    In addition we have:
    $\sup_{\lambda\ge0,\nu\in\mathbb{R}^{n_h}}q(\lambda,\nu)=\inf_{x\in\mathcal{C}}f(x)$
  \end{sstOnlyFrame}
\end{sstTitleBox}

\textbf{How to solve any KKT problem (exam tactic):}
\begin{enumerate}[leftmargin=*,itemsep=-1pt]
  \item Write problem as \(\min f(x)\) s.t. \(h(x) \leq 0, g(x) = 0\)
    (ensure \(\leq 0\) format).
  \item Write Lagrangian: \(\mathcal{L}(x, \lambda, \nu) = f(x) +
    \lambda h(x) + \nu g(x)\).
  \item List the 4 KKT conditions (as above).
  \item Solve by branching on complementary slackness:
    \begin{itemize}[leftmargin=*,itemsep=-1pt]
      \item \textit{Case A:} Assume \(\lambda = 0\) (constraint
        inactive), solve \(\nabla_x \mathcal{L}=0\) for \(x\), check
        if \(h(x)\leq 0\).
      \item \textit{Case B:} Assume \(\lambda > 0 \implies h(x)=0\)
        (constraint active), use this equality to solve for
        \(x,\lambda\), check \(\lambda \geq 0\).
    \end{itemize}
\end{enumerate}
\textbf{Tactic:} Write full \(\mathcal L\), set \(\nabla_x\mathcal
L=0\), substitute \(x^*\) into \(d(\lambda)\), maximize (often
quadratic). Check Slater. Compl.\ slackness key.

\textbf{Example} \(\min\frac12\|a\|^2\) s.t.\ \(y_i(a^\top
x_i+b)\ge1\); dual \(\max_{\lambda\ge0}d(\lambda)\). Solves to
\(\lambda_A^*=\lambda_B^*=1/4\).

\textbf{Example} \(f=1/x_1+4/x_2\), \(x_1+x_2\le3\),
\(x_{1,2}\ge1\): write full \(\mathcal L\), solve stationarity+slack.

\section*{ADMM Derivation (mechanical steps)}

\textbf{Step 1:} Reformulate problem to have separated variables:
\[
  \min_{x,z} f(x) + g(z) \quad \text{s.t. } Ax - z = c.
\]

\textbf{Step 2:} Write Augmented Lagrangian:
\[
  \mathcal{L}_\rho(x,z,\lambda) = f(x) + g(z) + \lambda^\top(Ax - z -
  c) + \frac{\rho}{2}\|Ax - z - c\|_2^2.
\]

\textbf{Step 3:} Update \(x\): Set \(\nabla_x \mathcal{L}_\rho = 0\).
This usually results in a linear system:
\[
  (A^\top A + \text{something})x = \dots
\]

\textbf{Step 4:} Update \(z\): If \(g(z)\) is an \(L_1\) norm (e.g.,
\(|z|_1\)), the solution is exactly the proximal operator/soft-thresholding:
\[
  z_{k+1} = \text{prox}_{1/\rho}(Ax_{k+1} - c + \lambda_k/\rho).
\]

\textbf{Step 5:} Update \(\lambda\):
\[
  \lambda_{k+1} = \lambda_k + \rho(Ax_{k+1} - z_{k+1} - c).
\]
\textbf{Tactic:} Memorize the 5 steps; the \(x\)-update reduces to
solving a linear system; the \(z\)-update is often a closed-form
proximal operator.

\section*{Reformulating max and absolute values (Epigraph)}

When minimizing \(\max(f(x), g(x))\) or \(|f(x)|\):

\begin{itemize}[leftmargin=*,itemsep=-1pt]
  \item Introduce an auxiliary variable \(t\).
  \item New objective: \(\min t\).
  \item For \(\max\): constraints are \(f(x) \le t\) and \(g(x

